% =====================================================================
%  Rozdział do monografii:
%  „Transformacja cyfrowa sektora publicznego – cyberbezpieczeństwo,
%   regulacje prawne i sztuczna inteligencja” (Wydawnictwo Naukowe UMK)
%
%  Kompilacja:  pdflatex rozdzial && biber rozdzial && pdflatex rozdzial
%               (lub: make)
%  Bibliografia: bibliografia.bib
% =====================================================================
\documentclass[12pt,a4paper]{article}

% ---------------------------------------------------------------------
%  Język, kodowanie, czcionka (Times New Roman → klon newtx)
% ---------------------------------------------------------------------
\usepackage[T1]{fontenc}
\usepackage[utf8]{inputenc}
\usepackage[polish]{babel}
\usepackage{newtxtext,newtxmath}
\usepackage[autostyle=true]{csquotes}   % „cytat” wg zasad polskich

% ---------------------------------------------------------------------
%  Układ strony: marginesy 2,5 cm, interlinia 1,5, wcięcie 0,75 cm,
%  wyjustowanie, ciągła numeracja stron
% ---------------------------------------------------------------------
\usepackage[margin=2.5cm]{geometry}
\usepackage{setspace}
\onehalfspacing
\setlength{\parindent}{0.75cm}
\setlength{\parskip}{0pt}
\usepackage{indentfirst}
\pagestyle{plain}

% ---------------------------------------------------------------------
%  Śródtytuły: 12 pt, pogrubienie, wcięcie 0,75 cm; sekcje numerowane,
%  podsekcje bez numeracji
% ---------------------------------------------------------------------
\usepackage{titlesec}
\setcounter{secnumdepth}{1}
\titleformat{\section}{\normalfont\normalsize\bfseries}{\thesection.}{0.5em}{}
\titlespacing*{\section}{0.75cm}{1.5ex plus .5ex minus .2ex}{0.5ex plus .2ex}
\titleformat{\subsection}{\normalfont\normalsize\bfseries\itshape}{}{0pt}{}
\titlespacing*{\subsection}{0.75cm}{1.2ex plus .5ex minus .2ex}{0.5ex plus .2ex}

% ---------------------------------------------------------------------
%  Przypisy: 10 pt, interlinia pojedyncza, wysunięcie 0,25 cm,
%  numeracja ciągła, odnośnik w indeksie górnym
% ---------------------------------------------------------------------
\makeatletter
\renewcommand\@makefntext[1]{%
  \setlength{\hangindent}{0.25cm}\hangafter=1
  \noindent\textsuperscript{\@thefnmark}\,#1}
\makeatother
\setlength{\skip\footins}{1.5em plus 4pt minus 2pt}

% ---------------------------------------------------------------------
%  Tabele i ilustracje: numer i tytuł / podpis
% ---------------------------------------------------------------------
\usepackage{graphicx}
\graphicspath{{rysunki/}}
\usepackage{booktabs}
\usepackage{tabularx}
\newcolumntype{L}{>{\raggedright\arraybackslash}X}
\newcolumntype{P}[1]{>{\raggedright\arraybackslash}p{#1}}
\usepackage{caption}
\captionsetup{font={small,singlespacing},labelfont=bf,labelsep=period,
  justification=centering,skip=6pt}
\captionsetup[table]{position=top}
% Źródło pod tabelą/rysunkiem, np. \zrodlo{opracowanie własne}
\newcommand{\zrodlo}[1]{\par\smallskip{\footnotesize Źródło: #1.}\par}

% ---------------------------------------------------------------------
%  Bibliografia: biblatex + biber, przypisy wg wzoru Wydawnictwa UMK
% ---------------------------------------------------------------------
\usepackage[
  backend=biber,
  style=verbose-trad2,     % pełny opis przy 1. powołaniu, Tamże, dz. cyt.
  sorting=nty,
  dashed=false,            % pełne nazwisko przy każdej pozycji
  giveninits=true,         % J. Nowak / Nowak J.
  maxnames=3, minnames=1,
  idemtracker=false,       % bez „tenże”
  loccittracker=false,     % bez „loc. cit.”
  dateabbrev=false,
  isbn=false, eprint=false,
  doi=true, url=true,
]{biblatex}
\addbibresource{bibliografia.bib}

% --- interpunkcja: przecinki między elementami opisu
\renewcommand*{\newunitpunct}{\addcomma\space}
\renewcommand*{\intitlepunct}{\addspace}
\DeclareDelimFormat{finalnamedelim}{\addcomma\space}
\DeclareDelimFormat[bib,biblist]{nametitledelim}{%
  \ifnameundef{author}{\addcolon\space}{\addcomma\space}}% „(red.): Tytuł”
\DeclareDelimFormat{editortypedelim}{\addspace}

% --- nazwiska: w przypisach „J. Nowak”, w bibliografii „Nowak J.”
\DeclareNameAlias{default}{given-family}
\DeclareNameAlias{sortname}{given-family}
\DeclareNameAlias{labelname}{given-family}
\renewcommand*{\revsdnamepunct}{}
\AtBeginBibliography{%
  \DeclareNameAlias{default}{family-given}%
  \DeclareNameAlias{sortname}{family-given}}

% --- (red.)
\DeclareFieldFormat{editortype}{\mkbibparens{#1}}

% --- tytuły kursywą, tytuł czasopisma w cudzysłowie, akty prawne prosto
\DeclareFieldFormat*{title}{\mkbibemph{#1}}
\DeclareFieldFormat*{booktitle}{\mkbibemph{#1}}
\DeclareFieldFormat*{citetitle}{\mkbibemph{#1}}
\DeclareFieldFormat{journaltitle}{\mkbibquote{#1}}
\DeclareFieldFormat[legislation]{title}{#1}
\DeclareFieldFormat[legislation]{citetitle}{#1}

% --- strony: „s. 5”; zakres stron rozdziału nie jest drukowany
%     (bibliografia „bez numerów stron”, w przypisie strony = postnote)
\DeclareFieldFormat{postnote}{\mkpageprefix[pagination]{#1}}
\AtEveryBibitem{\clearfield{pages}}
\AtEveryCitekey{\clearfield{pages}}

% --- „(w:)” + redaktorzy przed tytułem tomu zbiorowego
\renewbibmacro*{in:}{%
  \ifentrytype{article}
    {}
    {\setunit{\addspace}%
     \printtext{(w\addcolon)}%
     \setunit{\addspace}%
     \ifnameundef{editor}
       {}
       {\printnames{editor}%
        \setunit{\addspace}%
        \printtext{\mkbibparens{\bibstring{editor}}\addcolon}%
        \clearname{editor}%
        \setunit{\addspace}}}}

% --- czasopismo: „Tytuł czasopisma” 2000, t. 7, nr 1 (12), e1211
\renewbibmacro*{journal+issuetitle}{%
  \usebibmacro{journal}%
  \setunit*{\addspace}%
  \printfield{year}%
  \setunit{\addcomma\space}%
  \printfield{volume}%
  \setunit{\addcomma\space}%
  \printfield{number}%
  \setunit{\addcomma\space}%
  \printfield{eid}%
  \newunit}
\DeclareFieldFormat[article]{volume}{t\adddot~#1}
\DeclareFieldFormat[article]{number}{nr~#1}

% --- wydawnictwo, miejsce rok: „Wydawnictwo, Warszawa 2000”
\renewbibmacro*{publisher+location+date}{%
  \printlist{publisher}%
  \setunit*{\addcomma\space}%
  \printlist{location}%
  \setunit*{\addspace}%
  \usebibmacro{date}%
  \newunit}

% --- DOI jako pełny odnośnik, bez daty dostępu; w przypisie tylko przy
%     pierwszym powołaniu i po numerze strony: „…, s. 15, https://doi.org/…”
\newcommand*{\urlDOI}[1]{\url{https://doi.org/#1}}
\DeclareFieldFormat{doi}{\urlDOI{#1}}
\AtEveryBibitem{\iffieldundef{doi}{}{\clearfield{url}\clearfield{urlyear}}}
\AtEveryCitekey{%
  \global\let\cytDOI\relax
  \iffieldundef{doi}
    {}
    {\clearfield{url}\clearfield{urlyear}%
     \ifciteseen{}{\savefield{doi}{\cytDOItmp}\global\let\cytDOI\cytDOItmp}%
     \clearfield{doi}}}
\renewbibmacro*{cite:postnote}{%
  \usebibmacro{postnote}%
  \ifx\cytDOI\relax\else
    \setunit{\addcomma\space}%
    \printtext{\expandafter\urlDOI\expandafter{\cytDOI}}%
    \global\let\cytDOI\relax
  \fi}

% --- strony WWW: „www.nazwastrony.pl (dostęp 01.06.2026 r.)”
\DeclareFieldFormat{url}{\url{#1}}
\renewbibmacro*{url+urldate}{%
  \printfield{url}%
  \iffieldundef{urlyear}
    {}
    {\setunit*{\addspace}%
     \printtext[parens]{%
       \bibstring{urlseen}\addspace
       \mkdayzeros{\thefield{urlday}}.\mkmonthzeros{\thefield{urlmonth}}.%
       \thefield{urlyear}\addspace r\adddot}}}

% --- akty prawne: „Ustawa z dnia … (Dz.U. … z późn. zm.)”
\DeclareBibliographyDriver{legislation}{%
  \usebibmacro{bibindex}%
  \usebibmacro{begentry}%
  \printfield{title}%
  \setunit{\addspace}%
  \printtext[parens]{\printfield{note}}%
  \usebibmacro{finentry}}

% --- kolejne powołanie aktu prawnego: sam skrót (np. „RODO”), bez „dz. cyt.”
\renewbibmacro*{cite:title}{%
  \printtext[bibhyperlink]{%
    \printfield[citetitle]{labeltitle}%
    \ifentrytype{legislation}
      {}
      {\setunit{\printdelim{nametitledelim}}%
       \bibstring[\mkibid]{opcit}}}}

% --- bibliografia: 10 pt, wyjustowanie, wysunięcie 0,25 cm
\renewcommand*{\bibfont}{\footnotesize\setstretch{1}}
\setlength{\bibhang}{0.25cm}
\setlength{\bibitemsep}{0.3\baselineskip}
\defbibheading{bibliography}[Bibliografia]{\section*{#1}}

\usepackage[hidelinks]{hyperref}
\urlstyle{same}

% ---------------------------------------------------------------------
%  Pierwsza strona rozdziału
% ---------------------------------------------------------------------
\makeatletter
\newcommand{\autorzy}[1]{\def\@autorzy{#1}}
\newcommand{\afiliacja}[1]{\def\@afiliacja{#1}}
\newcommand{\tytulPL}[1]{\def\@tytulPL{#1}}
\newcommand{\tytulEN}[1]{\def\@tytulEN{#1}}
\newcommand{\naglowekRozdzialu}{%
  \begingroup
  \singlespacing
  \noindent{\bfseries\@autorzy}\par
  \noindent\@afiliacja\par
  \vspace{2\baselineskip}
  \begin{center}\bfseries\MakeUppercase{\@tytulPL}\\[0.5\baselineskip]
    \MakeUppercase{\@tytulEN}\end{center}
  \vspace{\baselineskip}
  \endgroup}
\makeatother

% Abstrakty: Times New Roman 10 pt
\newenvironment{streszczenie}[2]% #1 = etykieta abstraktu, #2 = etykieta słów kluczowych
  {\def\etykietaSlow{#2}\par\footnotesize\singlespacing\noindent\textbf{#1:}\space\ignorespaces}
  {\par}
\newcommand{\slowakluczowe}[1]{\par\smallskip\noindent\textbf{\etykietaSlow:} #1\par\medskip}

% =====================================================================
\begin{document}
% =====================================================================

\autorzy{Bartosz LEWANDOWSKI}
\afiliacja{Atende S.A.}
\tytulPL{Ochrona przed wyciekiem danych: rola systemów DLP w~wielowarstwowym modelu cyberbezpieczeństwa}
\tytulEN{Data leak protection: the role of DLP systems in a~multi-layered cybersecurity model}

\naglowekRozdzialu

\begin{streszczenie}{Abstrakt}{Słowa kluczowe}
Przedmiotem rozważań jest rola systemów zapobiegania wyciekom danych (DLP)
w~wielowarstwowym modelu cyberbezpieczeństwa. Rozważania mają charakter
przeglądowo-koncepcyjny i~opierają się na literaturze naukowej, 
przepisach RODO oraz raportach branżowych. Wyciek danych potraktowano
jako skutek łańcucha zdarzeń, w~którym błędy użytkowników nakładają się na
luki organizacyjne i~techniczne. Do opisu tego zjawiska wykorzystano model
sera szwajcarskiego J.~Reasona. Przedstawiono mechanizmy DLP: klasyfikację
informacji, ochronę danych w~trzech stanach, analizę treści i~kontekstu oraz
stopniowanie reakcji. Porównano zakres DLP z~systemami EDR, NDR i~SIEM.
Zaproponowano ujęcie DLP jako uzupełnienia triady widoczności SOC
o~perspektywę chronionych informacji. Omówiono warunki skutecznego wdrożenia
oraz finansowe, prawne i~społeczne konsekwencje wycieków. Analiza wskazuje,
że DLP skutecznie ogranicza wycieki przypadkowe, ale nie zastępuje innych
warstw ochrony i~wymaga wsparcia procedur organizacyjnych.
\slowakluczowe{zapobieganie wyciekom danych, DLP, poufność informacji,
wielowarstwowa ochrona, SOC, model sera szwajcarskiego, RODO}
\end{streszczenie}

% TODO: wersja robocza tłumaczenia – do weryfikacji przez autora
\begin{streszczenie}{Abstract}{Key words}
The article discusses the role of data loss prevention (DLP) systems in
a~multi-layered cybersecurity model. It is a~review and conceptual study
based on scientific literature,  the GDPR and industry
reports. A~data leak is treated as the outcome of a~chain of events in which
user errors coincide with organisational and technical gaps, described with
J.~Reason's Swiss cheese model. The article presents the mechanisms of DLP:
information classification, protection of data in three states, content and
context analysis, and graduated response. It compares the scope of DLP with
EDR, NDR and SIEM systems and proposes to view DLP as a~complement to the SOC
visibility triad, adding the perspective of protected information. The
conditions for effective deployment and the financial, legal and social
consequences of leaks are discussed. The analysis shows that DLP effectively
limits accidental leaks, but does not replace other layers of protection and
requires the support of organisational procedures.
\slowakluczowe{data loss prevention, DLP, information confidentiality,
defence in depth, SOC, Swiss cheese model, GDPR}
\end{streszczenie}

% ---------------------------------------------------------------------
\section{Wprowadzenie}

Poufność należy, obok integralności i~dostępności, do podstawowych
właściwości chronionych w~systemach bezpieczeństwa informacji. Jej naruszenie
trudno odwrócić. Zniszczone dane można odtworzyć z~kopii zapasowej, natomiast
danych raz ujawnionych nie da się ponownie uczynić poufnymi.

Organizacje przetwarzają informacje jednocześnie na stacjach roboczych,
w~usługach chmurowych, w~poczcie elektronicznej, w~narzędziach do współpracy
i~w~systemach partnerów. Każde z~tych miejsc może stać się kanałem wycieku.
Zapory sieciowe i~systemy kontroli dostępu rozstrzygają, kto i~skąd może się
połączyć. Nie oceniają jednak, czy konkretna treść powinna opuścić
organizację. Tę funkcję pełnią systemy zapobiegania wyciekom danych (ang.
\emph{Data Loss Prevention}, DLP). Są to narzędzia, które rozpoznają
informacje chronione w~treści i~kontekście operacji na danych, takich jak
zapis na nośniku, wysłanie wiadomości czy przesłanie pliku do usługi
chmurowej. Zależnie od przyjętej polityki rejestrują taką operację,
ostrzegają użytkownika albo ją blokują\footcite{alneyadi2016survey}.

Samo wdrożenie takiego narzędzia nie zapewnia jednak ochrony. DLP rozpoznaje
treść, ale nie zastępuje systemów obserwujących urządzenia, sieć i~dzienniki
zdarzeń. Jego skuteczność zależy od tego, czy organizacja wie, jakie
informacje chroni, i~czy potrafi utrzymać reguły bez mnożenia wyjątków.
Wyciek rzadko jest skutkiem pojedynczego błędu. Zwykle kończy łańcuch
zdarzeń, w~którym kolejne zabezpieczenia okazały się niewystarczające.
Dlatego DLP warto rozpatrywać jako jedną z~warstw wielowarstwowego modelu
cyberbezpieczeństwa, a~jego wartość oceniać na tle zagrożeń, które ma
ograniczać, i~kosztów, które pojawiają się, gdy ochrona zawiedzie. Punktem
wyjścia jest pytanie, dlaczego kanałów wycieku wciąż przybywa.

% ---------------------------------------------------------------------
\section{Geneza i~współczesne uwarunkowania wycieków danych}

Ryzyko wycieku rośnie wraz z~liczbą miejsc, w~których informacja jest
przechowywana i~przez które przepływa. Long Cheng, Fang Liu i~Danfeng Yao
wiążą je z~potrzebą gromadzenia i~analizowania coraz większych zbiorów danych
oraz z~intensywnym korzystaniem z~nowoczesnych kanałów
komunikacji\footcite{cheng2017enterprise}. Autorzy wymieniają wśród możliwych
wektorów wycieku m.in. usługi udostępniania plików w~chmurze, pocztę
elektroniczną, komunikatory, nośniki wymienne, podatności baz danych, kradzież
laptopów, utratę kopii zapasowych i~serwisy społecznościowe. Lista ta
pokazuje, że wyciek nie jest zjawiskiem jednego kanału.

Praca rozproszona przesunęła przetwarzanie danych poza siedzibę organizacji.
Ten sam przegląd wskazuje, że współdzielenie danych z~partnerami i~klientami,
współpraca przez chmurę oraz praca poza siedzibą firmy zwiększają
prawdopodobieństwo wycieku\footcite{cheng2017enterprise}. Nowsze dane
potwierdzają znaczenie podmiotów zewnętrznych. Według komunikatu Verizon
dotyczącego raportu \emph{Data Breach
Investigations Report} (DBIR) z~2026~roku, opartego na danych z~2025~roku, udział strony
trzeciej odnotowano w~48\% analizowanych naruszeń, a~wskaźnik ten wzrósł
o~60\% względem poprzedniej edycji\footcite{verizon2026dbir}. Granica
organizacji przestaje więc pokrywać się z~granicą jej sieci.

Narzędzia generatywnej sztucznej inteligencji tworzą kolejny kanał przepływu
informacji. Ten sam komunikat podaje, że korzystanie przez pracowników
z~niezatwierdzonych narzędzi sztucznej inteligencji (ang.
\emph{artificial intelligence}, AI), określane jako \emph{shadow AI}, stało się trzecią
najczęstszą niezłośliwą aktywnością związaną z~wyciekiem danych. Odsetek
pracowników często korzystających z~takich narzędzi wzrósł w~ciągu roku
z~15\% do 45\%\footcite{verizon2026dbir}. Z~perspektywy ochrony informacji
wklejenie dokumentu do zewnętrznej usługi jest przekazaniem danych poza
kontrolę organizacji, niezależnie od intencji użytkownika. Ta interpretacja
wyjaśnia, dlaczego problem dotyczy bardziej polityk ochrony danych niż samej
technologii AI.

Wyciek może mieć trzy różne źródła: działanie napastnika, celowe nadużycie
uprawnień przez osobę z~wewnątrz organizacji (ang. \emph{insider}), czyli
pracownika lub współpracownika mającego uprawniony dostęp do jej zasobów,
albo niezamierzony błąd. Cheng
i~współautorzy porządkują je według dwóch kryteriów: intencji (celowe lub
przypadkowe) oraz pochodzenia sprawcy (wewnętrzny lub
zewnętrzny)\footcite{cheng2017enterprise}. Do zagrożeń przypadkowych zaliczają
m.in. omyłkową publikację, błąd konfiguracji, niewłaściwe szyfrowanie i~utratę
komputera. Obie grupy mają polskie przykłady. W~2024~roku Prezes UODO nałożył
ponad 238~tys.~zł kary na spółkę, której pracownik zgubił pendrive
z~niezaszyfrowanymi danymi osobowymi innego pracownika\footcite{uodo2024pendrive}.
Organ wskazał, że administrator uwzględnił w~analizie ryzyka kradzież
i~zniszczenie nośnika, ale pominął jego zwykłą utratę. W~sprawie spółki
Fortum Marketing and Sales Polska źródłem naruszenia był błąd konfiguracji.
Podmiot przetwarzający utworzył podczas zmian w~systemie dodatkową bazę danych
klientów na serwerze bez odpowiednio skonfigurowanych zabezpieczeń. Osoby
nieuprawnione skopiowały tę bazę, a~administrator dowiedział się o~incydencie
od internautów\footcite{uodo2022fortum}.

Podział ten ma konsekwencje praktyczne. Napastnik zewnętrzny musi najpierw
uzyskać dostęp, osoba z~wewnątrz już go ma, a~pracownik popełniający błąd nie
ukrywa swoich działań. Każdy z~tych scenariuszy pozostawia inne ślady
i~wymaga innego sposobu wykrywania.

% ---------------------------------------------------------------------
\section{Charakterystyka zagrożeń w~świetle badań i~raportów}

Obraz zagrożeń opiera się w~dużej mierze na raportach branżowych, ponieważ
organizacje rzadko udostępniają dane o~własnych incydentach badaczom
akademickim. Dwie najczęściej przywoływane serie to DBIR firmy Verizon oraz \emph{Cost of a~Data Breach}
przygotowywany przez Ponemon Institute na zlecenie IBM. Obie serie mają
odmienną metodykę, więc ich wyniki należy czytać osobno.

Raport DBIR 2026 wskazuje wykorzystanie podatności jako najczęstszy punkt
wejścia. Odpowiada ono za 31\% naruszeń i~po raz pierwszy w~historii raportu
wyprzedziło użycie skradzionych poświadczeń\footcite{verizon2026dbir}. Ten sam
komunikat podaje, że socjotechnika prowadzona kanałami mobilnymi (wiadomości
tekstowe, połączenia głosowe) osiąga skuteczność o~40\% wyższą niż tradycyjny
phishing, czyli wyłudzanie danych w~wiadomościach e-mail podszywających się
pod zaufanego nadawcę. Dane te dotyczą głównie ataków zewnętrznych, ale pokazują
też, że użytkownik pozostaje częstym celem ataku.

Udział osób z~wewnątrz organizacji trudniej oszacować, ponieważ zależy on od
przyjętej definicji i~źródła danych. Pełna wersja raportu DBIR 2026 obejmuje
ponad 22~tys. potwierdzonych naruszeń z~okresu od listopada 2024 do
października 2025~roku. Sprawcy wewnętrzni wystąpili w~12\% z~nich, wobec
18\% w~poprzedniej edycji\footcite[22]{verizon2026dbirraport}. Autorzy
raportu zaznaczają, że większość tych przypadków to pomyłki, błędy
konfiguracji i~obchodzenie procedur, a~nie sabotaż. Umyślne nadużycie
uprawnień odpowiadało za niespełna 4\% naruszeń. Najczęstszym motywem
w~tej grupie była wygoda (60\%), np. przesłanie danych służbowych na
prywatną skrzynkę w~celu pracy z~domu\footcite[57--58]{verizon2026dbirraport}.
Inaczej wygląda administracja publiczna. Sprawcy wewnętrzni odpowiadali tam
za 44\% naruszeń, głównie z~powodu korespondencji doręczonej niewłaściwemu
adresatowi, która stanowiła 88\% błędów w~tym
sektorze\footcite[89, 92]{verizon2026dbirraport}.

Inną perspektywę daje badanie Ponemon Institute przeprowadzone na zlecenie
firmy DTEX Systems wśród 8750 specjalistów IT i~bezpieczeństwa z~354
organizacji, które doświadczyły incydentów wewnętrznych. W~2025~roku 53\%
tych incydentów wynikało z~zaniedbań pracowników lub kontraktorów, 27\%
z~działań celowych, a~20\% z~kradzieży poświadczeń. Średni roczny koszt
obsługi incydentów wewnętrznych wyniósł 19,5~mln USD na
organizację\footcite{ponemon2026insider}. Próba obejmuje wyłącznie
organizacje dotknięte takimi incydentami, więc wyniki nie opisują ogółu
przedsiębiorstw. Oba źródła, mimo odmiennej metodyki, prowadzą do zbieżnego
wniosku: po stronie wewnętrznej przeważa błąd lub zaniedbanie osoby
z~uprawnionym dostępem, a~nie działanie w~złej wierze. Przegląd
systematyczny 121 prac opublikowanych w~latach 2014--2025 wskazuje
jednocześnie na lukę badawczą. Niewiele metod wykrywania zagrożeń
wewnętrznych uwzględnia czynniki ludzkie, przede wszystkim z~powodu braku
zbiorów danych łączących informacje o~zachowaniu i~kontekście użytkownika.
Autorzy przeglądu podkreślają potrzebę takich modeli zwłaszcza dla zagrożeń
nieumyślnych\footcite{pathirana2026human}.

Czynniki ludzkie, techniczne i~organizacyjne występują zwykle łącznie.
W~ujęciu Chenga i~współautorów przypadkowe wycieki wynikają głównie ze słabo
zaprojektowanych procesów biznesowych, czyli z~braku odpowiednich technologii
zapobiegawczych i~polityk bezpieczeństwa, albo z~przeoczenia ze strony
pracownika\footcite{cheng2017enterprise}. Błąd użytkownika jest tu więc
objawem, a~nie jedyną przyczyną. Podobnie można odczytać dane IBM: tylko 37\%
organizacji dotkniętych naruszeniem zadeklarowało szyfrowanie danych
wrażliwych zarówno w~spoczynku, jak i~podczas
przesyłania\footcite{ibm2026cost}.

Typowe scenariusze wycieku obejmują wyprowadzenie danych przez złośliwe
oprogramowanie po włamaniu, nadużycie uprawnień lub sabotaż ze strony
insidera, błędną konfigurację zasobu, omyłkową publikację oraz utratę
urządzenia z~niezaszyfrowanymi danymi\footcite{cheng2017enterprise}. Kanałem
może być sieć, punkt końcowy, usługa chmurowa albo nośnik fizyczny. Asaf
Shabtai, Yuval Elovici i~Lior Rokach traktują kanał wycieku jako jedną z~osi
klasyfikacji rozwiązań DLP, obok źródła wycieku i~stanu
danych\footcite{shabtai2012survey}.

\subsection{Ograniczenia interpretacji statystyk}

Statystyki naruszeń opisują różne zjawiska, choć często używają podobnych
słów. DBIR odróżnia incydent, czyli zdarzenie naruszające poufność,
integralność lub dostępność zasobu informacyjnego, od naruszenia, czyli
incydentu zakończonego potwierdzonym ujawnieniem danych osobie
nieuprawnionej\footcite{verizon2024dbir}. Samo potencjalne narażenie danych
nie spełnia tej definicji. Rozproszony atak odmowy usługi (ang.
\emph{Distributed Denial of Service}, DDoS), który przeciąża system
i~uniemożliwia korzystanie z~niego, jest więc zwykle incydentem, ale nie
naruszeniem.

Ogólne rozporządzenie o~ochronie danych (RODO), czyli unijny akt prawny
regulujący przetwarzanie danych osobowych, przyjmuje szersze pojęcie. Naruszenie ochrony danych osobowych obejmuje
nie tylko nieuprawnione ujawnienie lub dostęp, ale także przypadkowe lub
niezgodne z~prawem zniszczenie, utracenie i~zmodyfikowanie danych (art.~4
pkt~12)\footcite{rodo2016}. Zaszyfrowanie danych przez oprogramowanie
wymuszające okup (ang. \emph{ransomware}) bez ich wyprowadzenia może być naruszeniem w~rozumieniu RODO, ale
nie naruszeniem w~rozumieniu DBIR. Wyciek rozumiany jest tu węziej, jako
przekroczenie przez informację granicy, w~której jej przetwarzanie
jest uprawnione.

Różnice dotyczą też prób badawczych. DBIR opiera się na danych z~dochodzeń
prowadzonych przez Verizon i~organizacje
współpracujące\footcite{verizon2024dbir}, a~nie na losowej próbie
organizacji. Edycja raportu IBM z~2026~roku obejmuje 602 organizacje, które
doświadczyły naruszenia od marca 2025 do lutego 2026~roku\footcite{ibm2026cost}.
Odsetki z~tych raportów opisują więc zbiór zbadanych naruszeń lub
organizacji. Nie mówią o~tym, jak często naruszenia dotykają ogółu
przedsiębiorstw.

% ---------------------------------------------------------------------
\section{Istota i~zakres ochrony DLP}

Termin DLP nie ma jednej, powszechnie przyjętej definicji. Sultan Alneyadi,
Elankayer Sithirasenan i~Vallipuram Muthukkumarasamy określają systemy DLP
jako wyspecjalizowane mechanizmy wykrywania i~zapobiegania wyciekom danych
poufnych w~trakcie ich używania, przesyłania i~przechowywania\footcite{alneyadi2016survey}.
Zwracają jednocześnie uwagę, że samo pojęcie pozostaje niejednoznaczne, mimo
rosnącej liczby produktów komercyjnych. Cheng i~współautorzy używają
szerszego określenia DLPD (ang. \emph{Data Leak Prevention and Detection})
i~przypisują tym systemom trzy zadania: identyfikację, monitorowanie
i~zapobieganie ujawnieniu informacji wrażliwych, zarówno celowemu, jak
i~przypadkowemu\footcite{cheng2017enterprise}.

DLP różni się od podstawowych zabezpieczeń przedmiotem kontroli. Zapory
sieciowe, oprogramowanie antywirusowe, systemy wykrywania włamań, kontrola
dostępu i~szyfrowanie chronią systemy i~kanały. DLP ocenia treść lub kontekst
samych danych\footcite{cheng2017enterprise}. Cheng i~współautorzy formułują
też wymagania wobec praktycznego systemu DLP. Powinien blokować tylko
przepływy zawierające informacje wrażliwe i~przepuszczać zwykły ruch.
Powinien chronić przed nieostrożnymi pracownikami i~złośliwymi insiderami.
Powinien też utrudnić wyprowadzenie danych nawet wtedy, gdy zawiodą inne
zabezpieczenia.

\subsection{Identyfikacja i~klasyfikacja informacji}

System DLP może chronić tylko te informacje, które organizacja wcześniej
rozpoznała i~opisała. Ochrona poufności zaczyna się więc nie od wyboru
narzędzia, lecz od ustalenia, jakie dane organizacja przetwarza, gdzie się one
znajdują, kto z~nich korzysta i~jakie byłyby skutki ich ujawnienia. Na tej
podstawie dane przypisuje się do kategorii, na przykład publicznych,
wewnętrznych, poufnych i~ściśle chronionych, a~każdej kategorii
przyporządkowuje zasady przetwarzania. Dopiero tak opisane zasoby można
przełożyć na reguły systemu DLP. Klasyfikacja nie jest przy tym czynnością
jednorazową, ponieważ powstają nowe zbiory danych, systemy i~kanały ich
wymiany. Wyznacza ona zakres działania DLP. Błędy na tym etapie przenoszą się
na wszystkie dalsze decyzje systemu: dane pominięte w~klasyfikacji pozostają
bez ochrony, a~dane sklasyfikowane zbyt szeroko generują nadmiar alarmów.

\subsection{Dane przechowywane, przesyłane i~wykorzystywane}

Klasyczny podział obejmuje trzy stany danych. Cheng i~współautorzy
przypisują każdemu z~nich inny cel ochrony\footcite{cheng2017enterprise}.
Skanowanie danych przechowywanych na serwerach pozwala znaleźć zasoby
narażone na wyciek wewnątrz organizacji. Monitorowanie danych używanych na
stacjach roboczych ma zapobiegać niewłaściwemu obchodzeniu się z~nimi przez
użytkowników. Monitorowanie ruchu sieciowego ma powstrzymać przesłanie danych
poufnych poza sieć firmową. Shabtai i~współautorzy włączają stan danych do
taksonomii rozwiązań, obok źródła wycieku, kanału, sposobu wdrożenia,
charakteru działania (zapobiegawczy lub wykrywający) i~reakcji na
wyciek\footcite{shabtai2012survey}.

\subsection{Analiza treści i~kontekstu}

Systemy DLP rozpoznają dane chronione na dwa sposoby: przez analizę treści
albo przez analizę kontekstu\footcite{cheng2017enterprise}. Pierwsze podejście
sprawdza, co zawiera przesyłany lub zapisywany obiekt. Drugie bada metadane
i~okoliczności operacji: kto ją wykonuje, skąd, dokąd, w~jakim czasie i~czy
odbiega ona od typowego wzorca zachowania. Część rozwiązań łączy obie metody.
Tabela~\ref{tab:techniki} zestawia główne techniki wraz z~ich ograniczeniami.

\begin{table}[htbp]
  \centering
  \caption{Techniki wykrywania w~systemach DLP}
  \label{tab:techniki}
  \footnotesize\setstretch{1}
  \renewcommand{\arraystretch}{1.2}
  \begin{tabularx}{\textwidth}{@{}P{3.2cm}P{2cm}LL@{}}
    \toprule
    Technika & Rodzaj analizy & Zaleta & Ograniczenie \\
    \midrule
    Odciski cyfrowe (\emph{fingerprinting}) & treść
      & prosta implementacja, wykrywa całe dokumenty chronione
      & duża wrażliwość na modyfikację danych \\
    Wyrażenia regularne & treść
      & dobre dla danych o~przewidywalnym formacie (np. numery kart)
      & ograniczona ochrona, wysoki odsetek fałszywych alarmów \\
    Porównywanie zbiorów n-gramów & treść
      & toleruje drobne zmiany tekstu
      & duży koszt obliczeń i~pamięci \\
    Uczenie maszynowe & treść lub kontekst
      & odporność na modyfikacje danych
      & potrzeba dużych zbiorów uczących, złożoność \\
    Analiza zachowań & kontekst
      & ogranicza zagrożenia wewnętrzne
      & potrzeba dużych zbiorów uczących, wysoki odsetek fałszywych alarmów \\
    Znakowanie wodne & kontekst
      & wspiera analizę powykonawczą
      & podatność na usunięcie lub zniekształcenie znaku \\
    \bottomrule
  \end{tabularx}
  \zrodlo{opracowanie własne na podstawie: \citeauthor{cheng2017enterprise},
    \citetitle{cheng2017enterprise}}
\end{table}

Analiza treści skutecznie chroni przed przypadkową utratą danych, ale
świadomy sprawca może ją ominąć przez zaciemnienie
danych\footcite{cheng2017enterprise}. Ruch szyfrowany czyni ją bezużyteczną,
jeżeli system nie ma wglądu w~treść przed zaszyfrowaniem. Cheng
i~współautorzy przywołują też ocenę trzech rozwiązań komercyjnych
przeprowadzoną przez Davida Gugelmanna i~współautorów. Systemy te zapobiegały
przypadkowym wyciekom w~ruchu WWW, ale insider lub złośliwe oprogramowanie
mogły je łatwo obejść. Wynik ten wyznacza realistyczny zakres oczekiwań wobec
DLP.

\subsection{Monitorowanie, ostrzeganie i~blokowanie}

Reakcja na wykrytą operację może mieć różną siłę. System może jedynie
zarejestrować zdarzenie, ostrzec użytkownika i~poprosić o~uzasadnienie albo
zablokować operację. Shabtai i~współautorzy traktują działanie podejmowane po
wykryciu wycieku jako odrębną oś klasyfikacji
rozwiązań\footcite{shabtai2012survey}. Dobór reakcji jest decyzją
organizacyjną, a~nie wyłącznie techniczną. Blokada chroni dane, ale przy
błędnej regule zatrzymuje też uprawnioną pracę.

% ---------------------------------------------------------------------
\section{Miejsce DLP w~architekturze cyberbezpieczeństwa}

DLP nie zastępuje innych narzędzi wykrywania zagrożeń, ponieważ obserwuje
inny obiekt. Systemy EDR (ang. \emph{Endpoint Detection and Response})
analizują telemetrię punktów końcowych: uruchamiane procesy, zmiany
w~systemie i~aktywność na urządzeniu. Systemy NDR (ang. \emph{Network
Detection and Response}) analizują ruch sieciowy. Systemy SIEM (ang.
\emph{Security Information and Event Management}) gromadzą i~korelują
dzienniki zdarzeń z~wielu źródeł. DLP ocenia natomiast, czy dana operacja
dotyczy informacji chronionej i~czy jest zgodna z~polityką jej przetwarzania.
Tabela~\ref{tab:porownanie} zestawia te różnice.

\begin{table}[htbp]
  \centering
  \caption{Porównanie zakresów działania DLP, EDR, NDR i~SIEM}
  \label{tab:porownanie}
  \footnotesize\setstretch{1}
  \renewcommand{\arraystretch}{1.2}
  \begin{tabularx}{\textwidth}{@{}P{1.9cm}LLL@{}}
    \toprule
    Klasa rozwiązań & Przedmiot obserwacji & Pytanie, na które odpowiada
      & Typowa reakcja \\
    \midrule
    DLP & treść i~kontekst operacji na danych
      & czy ta informacja może opuścić dozwolony obszar?
      & rejestracja, ostrzeżenie użytkownika, blokada operacji \\
    EDR & procesy i~zmiany na punkcie końcowym
      & co dzieje się na urządzeniu?
      & izolacja hosta, zatrzymanie procesu \\
    NDR & ruch sieciowy
      & kto, z~kim i~w~jaki sposób się komunikuje?
      & alarm, wsparcie blokady połączenia \\
    SIEM & dzienniki zdarzeń z~wielu systemów
      & jak zdarzenia z~różnych źródeł łączą się w~całość?
      & korelacja, alarm, przekazanie do obsługi incydentu \\
    \bottomrule
  \end{tabularx}
  \zrodlo{opracowanie własne}
\end{table}

Koncepcję trzech źródeł widoczności w~centrum operacji bezpieczeństwa (ang.
\emph{Security Operations Center}, SOC), czyli zespole stale monitorującym
zdarzenia i~obsługującym incydenty, zaproponował w~2015~roku Anton Chuvakin, pracując wtedy jako analityk
Gartnera\footcite{chuvakin2020triad}. Pierwotnie nazwał ją \enquote{jądrową
triadą SOC}, przez analogię do triady środków przenoszenia broni jądrowej.
Trzy filary to dzienniki zdarzeń (SIEM), dane sieciowe i~telemetria punktów
końcowych (EDR). Dane sieciowe dostarczały pierwotnie narzędzia analizy ruchu
(ang. \emph{Network Traffic Analysis}, NTA), później określane jako NDR. Chuvakin zauważył później, że zakresy
tych źródeł częściowo się pokrywają. Dopuścił też, że widoczność aplikacji
może stać się czwartym filarem.

Triada odpowiada na pytanie, co dzieje się w~infrastrukturze. Nie odpowiada
wprost na pytanie, których informacji dotyczą obserwowane zdarzenia. Autor
proponuje traktować DLP jako uzupełnienie triady o~perspektywę chronionych
danych. Nie jest to stanowisko Gartnera ani Chuvakina, lecz interpretacja
proponowana w~niniejszym tekście. Przykładowo EDR może odnotować uruchomienie programu
archiwizującego, a~NDR wysyłkę dużego pliku do zewnętrznej usługi. Dopiero
DLP może wskazać, że archiwum zawiera dokumenty oznaczone jako poufne.

Wartość tej perspektywy rośnie, gdy zdarzenia DLP trafiają do wspólnego
procesu monitorowania. Cheng i~współautorzy pokazują wiele punktów, w~których
można rozmieścić uzupełniające się techniki ochrony: na urządzeniach
pracowników, serwerach pocztowych i~WWW, styku z~siecią zewnętrzną oraz
w~kontaktach z~dostawcami i~partnerami\footcite{cheng2017enterprise}.
Przewodnik SP 1800-28, opublikowany przez amerykański National Institute of
Standards and Technology (NIST), również opisuje rozwiązanie przykładowe złożone
z~wielu współpracujących systemów. Wykrywanie, reagowanie i~odtwarzanie po
naruszeniu poufności omówiono w~nim jako osobny projekt (SP
1800-29)\footcite{nist2024dataconf}.

Integracja ma znaczenie praktyczne dla obsługi incydentów. Alarm DLP bez
kontekstu z~EDR lub SIEM trudno ocenić: nie wiadomo, czy użytkownik działał
świadomie, czy jego konto zostało przejęte. Z~kolei alarm EDR bez informacji
o~rodzaju danych utrudnia ocenę skutków incydentu. Połączenie obu perspektyw
pozwala szybciej ustalić, czy doszło do wycieku i~kogo dotyczy.

% ---------------------------------------------------------------------
\section{Od pojedynczego błędu do wycieku: kontekst i~zachowania}

Wyciek rzadko wynika z~jednego działania. Dobrze to pokazuje analiza
naruszenia w~sieci handlowej Target z~2013~roku, przedstawiona przez Chenga
i~współautorów\footcite{cheng2017enterprise}. Napastnicy przejęli najpierw
przez phishing system dostawcy usług technicznych, który miał zdalny dostęp
do sieci Target. Następnie weszli do sieci firmy, zainstalowali złośliwe
oprogramowanie na terminalach płatniczych, przenieśli zebrane dane na
przejęte hosty wewnętrzne i~wyprowadzili je na zewnątrz. Każdy z~tych etapów
był osobną okazją do przerwania ataku.

Autorzy wskazują cztery techniczne przyczyny tego naruszenia. Target nie
stosował odpowiedniej kontroli dostępu wobec partnerów zewnętrznych. Nie
odseparował systemu płatniczego od reszty sieci. Nie zabezpieczył terminali
przed instalacją nieautoryzowanego oprogramowania. Miał wdrożone zapory
i~system zapobiegania włamaniom, ale nie zbadał ostrzeżeń, które te narzędzia
wygenerowały\footcite{cheng2017enterprise}. Ostatnia przyczyna jest
z~perspektywy DLP najważniejsza. Zabezpieczenie zadziałało, a~mimo to
bariera zawiodła na poziomie organizacji.

\subsection{Model sera szwajcarskiego}

Taką sekwencję dobrze opisuje model sera szwajcarskiego Jamesa
Reasona\footcite{reason2000human}. Reason przedstawia zabezpieczenia systemu
jako kolejne warstwy. Część z~nich jest techniczna (alarmy, bariery fizyczne,
automatyczne wyłączenia), część opiera się na ludziach, a~część na
procedurach i~kontroli administracyjnej. Każda warstwa ma luki, które
otwierają się, zamykają i~przesuwają. Luka w~pojedynczej warstwie zwykle nie
prowadzi do szkody. Szkoda powstaje dopiero wtedy, gdy luki w~wielu
warstwach chwilowo się pokryją.

Reason wyróżnia dwa źródła luk. Błędy aktywne to niebezpieczne działania
osób mających bezpośredni kontakt z~systemem: pomyłki, przeoczenia
i~naruszenia procedur. Warunki ukryte wynikają z~decyzji projektantów,
autorów procedur i~kierownictwa. Mogą pozostawać niezauważone przez lata,
zanim w~połączeniu z~błędem aktywnym umożliwią
wypadek\footcite{reason2000human}. W~odróżnieniu od błędów aktywnych warunki
ukryte można wykryć i~usunąć, zanim dojdzie do zdarzenia.

Analogia z~bezpieczeństwem lotniczym ma tu podstawę w~samym tekście Reasona.
Autor powołuje się na dane z~obsługi technicznej statków powietrznych, według
których około 90\% uchybień jakościowych uznano za niezawinione. Jako
przykłady organizacji wysokiej niezawodności wymienia m.in. ośrodki kontroli
ruchu lotniczego\footcite{reason2000human}. Reason podkreśla też, że skuteczne
zarządzanie ryzykiem wymaga kultury zgłaszania błędów, a~ta z~kolei wymaga
zaufania. Wniosek ten przenosi się na ochronę danych: pracownik, który obawia
się kary, rzadziej zgłosi omyłkowe wysłanie dokumentu.

W~ochronie przed wyciekiem danych błędem aktywnym może być wysłanie pliku do
niewłaściwego odbiorcy. Warunkami ukrytymi mogą być nadmierne uprawnienia,
brak klasyfikacji danych, nieodseparowane segmenty sieci albo brak procedury
analizy alarmów. Przypadek Target pokazuje, że w~praktyce występują one
łącznie\footcite{cheng2017enterprise}. DLP jest w~tym ujęciu jedną z~warstw.
Zmniejsza prawdopodobieństwo, że luki się pokryją, ale nie eliminuje luk
w~innych warstwach.

\subsection{Analiza kontekstu i~zachowań}

Reguła oceniająca pojedynczą operację odpowiada na pytanie, czy ta operacja
narusza politykę. Nie widzi jednak sekwencji. Rozważmy przykład hipotetyczny:
pracownik pobiera w~ciągu tygodnia więcej dokumentów niż zwykle, łączy je
w~archiwa i~przesyła do prywatnej usługi chmurowej. Każda z~tych czynności
osobno może mieścić się w~granicach uprawnień. Ryzyko ujawnia dopiero ich
układ.

Analiza zachowań próbuje uchwycić takie układy przez modelowanie typowej
aktywności użytkownika i~wykrywanie odchyleń\footcite{cheng2017enterprise}.
Cheng i~współautorzy zaznaczają jednak dwa ograniczenia. Pierwszym jest
wysoki odsetek fałszywych alarmów, wynikający ze zmienności kontekstu
w~czasie. Drugim jest luka semantyczna między obserwowalnymi zdarzeniami
w~systemie a~intencjami użytkownika, których nie da się zmierzyć
bezpośrednio. Insider działający w~granicach swoich uprawnień pozostaje więc
trudny do wykrycia.

Reakcja na pojedynczą operację i~analiza sekwencji pełnią różne funkcje.
Pierwsza zatrzymuje oczywiste naruszenia w~chwili ich wystąpienia. Druga
wskazuje sytuacje wymagające oceny przez analityka. Autor uważa, że łączenie
obu poziomów jest właściwsze niż próba zastąpienia jednego drugim. Blokowanie
na podstawie samej analizy zachowań, przy obecnym poziomie fałszywych
alarmów, zatrzymywałoby zbyt wiele uprawnionych działań.

% ---------------------------------------------------------------------
\section{Uwarunkowania wdrożenia i~skuteczności DLP}

Skuteczność DLP zależy bardziej od przygotowania organizacji niż od wyboru produktu. 
Równie istotnym czynnikiem sukcesu jest determinacja organizacji oraz odporność na 
wprowadzanie wyjątków do reguł. Kluczowe etapy wdrożeń można określić w następujący sposób.

\subsection{Zakres ochrony i~polityki}

Wdrożenie powinno zaczynać się od określenia, które informacje wymagają
ochrony i~dlaczego. Kolejność przyjęta w procesie to najpierw identyfikacja zasobów i~ryzyk, potem ochrona.
Organizacja, która nie wie, gdzie przechowuje dane osobowe klientów lub dokumentację projektową, nie napisze
dla nich trafnych reguł.

Polityki DLP powinny wynikać z~klasyfikacji danych i~z~procesów biznesowych.
Ta sama treść może być przesyłana zasadnie w~jednym procesie
i~niedopuszczalnie w~innym. Umowa z~kontrahentem wysłana do jego działu
prawnego jest operacją rutynową. Ta sama umowa wysłana na prywatny adres
pracownika wymaga wyjaśnienia. Reguła oparta wyłącznie na treści nie odróżni
tych przypadków. Potrzebny jest kontekst procesu.

\subsection{Etapowe uruchamianie}

Wdrożenie można podzielić na trzy etapy:
\begin{enumerate}
  \item Obserwacja. System rejestruje zdarzenia bez ingerencji w~pracę
    użytkowników. Organizacja poznaje rzeczywiste przepływy danych, które
    często różnią się od opisanych w~procedurach.
  \item Dostrajanie. Reguły są korygowane na podstawie zebranych zdarzeń.
    Użytkownicy otrzymują ostrzeżenia i~mogą podać uzasadnienie operacji.
  \item Blokowanie. Blokady obejmują najpierw reguły o~niskim odsetku
    fałszywych alarmów i~dane o~najwyższej wrażliwości, a~dopiero później
    kolejne kategorie.
\end{enumerate}

Kolejność ta wynika z~ograniczeń technik wykrywania. Wyrażenia regularne
i~analiza zachowań generują dużo fałszywych
alarmów\footcite{cheng2017enterprise}. Natychmiastowe włączenie blokad
zatrzymałoby więc część uprawnionej pracy. Użytkownicy zaczęliby wtedy
szukać kanałów omijających kontrolę.

\subsection{Fałszywe alarmy i~wyjątki}

Każdy fałszywy alarm kosztuje czas analityka i~obniża zaufanie do systemu.
Organizacja potrzebuje procedury rozpatrywania alarmów i~przyznawania
wyjątków. Wyjątek powinien mieć właściciela, uzasadnienie i~termin ważności.
Wyjątki bezterminowe stają się z~czasem warunkami ukrytymi w~rozumieniu
Reasona\footcite{reason2000human}: luką, o~której wszyscy zapomnieli.

Przypadek Target pokazuje drugą stronę tego problemu. Ostrzeżenia
wygenerowane przez narzędzia bezpieczeństwa nie zostały
zbadane\footcite{cheng2017enterprise}. System, który generuje alarmy, ale nie
ma przypisanej odpowiedzialności za ich obsługę, daje poczucie ochrony bez
rzeczywistej ochrony.

\subsection{Odpowiedzialność i~edukacja}

DLP wymaga podziału ról. Właściciele danych określają, co jest chronione.
Zespół bezpieczeństwa utrzymuje reguły i~obsługuje alarmy. Kierownicy
jednostek rozstrzygają o~wyjątkach w~swoich procesach. Dział prawny lub
inspektor ochrony danych ocenia, czy zdarzenie jest naruszeniem wymagającym
zgłoszenia.

Edukacja użytkowników ma w~tym modelu konkretny cel. Cheng i~współautorzy
wskazują, że zapobieganie przypadkowym wyciekom wymaga, obok środków
technicznych, budowania świadomości bezpieczeństwa w~miejscu
pracy\footcite{cheng2017enterprise}. Ostrzeżenie wyświetlane w~chwili
operacji może pełnić funkcję edukacyjną, jeżeli wyjaśnia, dlaczego operacja
została zatrzymana. Wsparciem jest kultura zgłaszania błędów opisana przez
Reasona\footcite{reason2000human}.

\subsection{Ocena skuteczności i~ograniczenia modelu}

Skuteczności DLP nie można mierzyć samą liczbą zablokowanych operacji. Duża
liczba blokad może oznaczać skuteczną ochronę albo źle dobrane reguły. Autor
proponuje kilka miar, które należy ocenić w~przyszłych badaniach: odsetek
alarmów uznanych po analizie za zasadne, czas od alarmu do decyzji
analityka, udział sklasyfikowanych repozytoriów objętych monitorowaniem oraz
liczbę aktywnych wyjątków bez wyznaczonego terminu.

Model ochrony oparty na DLP ma też granice. Analiza treści nie radzi sobie
z~danymi zaszyfrowanymi lub celowo zaciemnionymi. Insider działający
w~granicach swoich uprawnień może pozostać niewykryty. Brakuje też wspólnych
zbiorów danych testowych, co utrudnia porównywanie metod
wykrywania\footcite{cheng2017enterprise}. DLP ogranicza więc ryzyko wycieku,
ale go nie usuwa.

% ---------------------------------------------------------------------
\section{Finansowe i~społeczne konsekwencje wycieków}

Koszt naruszenia obejmuje znacznie więcej niż usunięcie jego technicznej
przyczyny. Według raportu IBM z~2026~roku średni globalny koszt naruszenia
wyniósł 4,99~mln USD\footcite{ibm2026cost}. Wartość ta jest średnią dla 602
organizacji, które doświadczyły naruszenia od marca 2025 do lutego
2026~roku. Nie opisuje typowego kosztu w~małej firmie ani w~polskich realiach
rynkowych. Średnia arytmetyczna jest też wrażliwa na pojedyncze bardzo
kosztowne przypadki.

Na koszty składają się obsługa incydentu, analiza powykonawcza, przywracanie
działalności, powiadomienia, wsparcie dla poszkodowanych i~utracone
przychody. Cheng i~współautorzy przytaczają dwa przykłady skali takich
skutków\footcite{cheng2017enterprise}. Target podał straty związane
z~naruszeniem z~2013~roku na poziomie 248~mln USD. Po ujawnieniu naruszeń
w~Yahoo Verizon zapłacił za przejęcie tej firmy o~350~mln USD mniej, niż
pierwotnie planowano. Drugi przykład pokazuje, że naruszenie może obniżyć
wartość przedsiębiorstwa w~transakcji kapitałowej.

\subsection{Konsekwencje prawne, kontraktowe i~reputacyjne}

W~Unii Europejskiej wyciek danych osobowych uruchamia obowiązki określone
w~RODO\footcite{rodo2016}. Administrator zgłasza naruszenie organowi
nadzorczemu bez zbędnej zwłoki, w~miarę możliwości nie później niż w~ciągu
72~godzin od jego stwierdzenia (art.~33). Jeżeli naruszenie może powodować
wysokie ryzyko dla praw lub wolności osób fizycznych, administrator
zawiadamia również te osoby (art.~34). Osoba, która poniosła szkodę
majątkową lub niemajątkową, ma prawo do odszkodowania (art.~82). Za
naruszenie przepisów rozporządzenia grozi administracyjna kara pieniężna,
w~najwyższym progu do 20~mln EUR lub do 4\% całkowitego rocznego światowego
obrotu przedsiębiorstwa (art.~83).

Obowiązek zgłoszenia w~ciągu 72~godzin łączy się bezpośrednio z~rolą
DLP. Administrator musi w~krótkim czasie ustalić, jakie dane zostały
ujawnione i~kogo dotyczą. Zdarzenia rejestrowane przez DLP, powiązane
z~klasyfikacją danych, mogą skrócić tę analizę. Jest to wniosek autora, a~nie
wymóg prawny.

Skutki kontraktowe i~reputacyjne trudniej zmierzyć. Wyciek danych
powierzonych przez kontrahenta może oznaczać naruszenie umowy i~utratę
klienta. Raport IBM wskazuje, że sprawcy ataków ransomware coraz częściej
wywierają presję przez zagrożenie dla reputacji. Najczęściej wykorzystują
reputację marki ofiary (41\% przypadków), dane pracowników (35\%)
i~własność intelektualną (31\%)\footcite{ibm2026cost}. Reputacja staje się
więc bezpośrednim narzędziem wymuszenia.

\subsection{Utrata zaufania i~skutki dla osób, których dane ujawniono}

Wyciek narusza zaufanie klientów, pracowników i~partnerów. Zaufanie
odbudowuje się wolniej niż systemy informatyczne, a~jego utraty nie oddają
w~pełni ani koszty obsługi incydentu, ani wysokość kar.

Najpoważniejsze skutki ponoszą osoby, których dane ujawniono. W~incydencie
australijskiego Red Cross Blood Service ujawnione dokumenty zawierały nie
tylko imiona, adresy i~daty urodzenia dawców. Obejmowały też informacje
o~aktywności seksualnej, przyjmowaniu narkotyków i~historii
chorób\footcite{cheng2017enterprise}. Takiej szkody nie da się naprawić
zmianą hasła ani wydaniem nowej karty płatniczej. Ujawnienie danych
identyfikacyjnych zwiększa z~kolei ryzyko kradzieży tożsamości, która była
typem szkody w~części największych naruszeń zestawionych przez tych samych
autorów.

Ochrona informacji jest więc odpowiedzialnością organizacyjną i~społeczną.
Organizacja przetwarza dane osób, które zwykle nie mają wpływu na sposób ich
zabezpieczenia. Skutki błędów ponoszą w~dużej mierze ci, którzy ich nie
popełnili.

% ---------------------------------------------------------------------
\section{Podsumowanie i~wnioski}

Systemy DLP wnoszą do wielowarstwowej ochrony informacji perspektywę, której
nie zapewniają narzędzia triady widoczności SOC. EDR, NDR i~SIEM opisują
zdarzenia na urządzeniach, w~sieci i~w~dziennikach systemów. DLP wiąże te
zdarzenia z~konkretnymi informacjami i~z~zasadami ich przetwarzania. Ta
zdolność jest użyteczna zarówno w~zapobieganiu wyciekom, jak i~w~ocenie
skutków incydentu.

Przegląd literatury wskazuje jednocześnie wyraźne granice tej klasy
rozwiązań. DLP skutecznie ogranicza wycieki przypadkowe, ale świadomy
sprawca może obejść analizę treści przez zaciemnienie lub szyfrowanie
danych\footcite{cheng2017enterprise}. Analiza zachowań rozszerza możliwości
wykrywania zagrożeń wewnętrznych, lecz generuje wiele fałszywych alarmów. DLP
należy więc traktować jako jedną z~warstw ochrony, a~nie jako rozwiązanie
samodzielne.

Model sera szwajcarskiego\footcite{reason2000human} pozwala ująć wyciek jako
skutek nakładania się błędów aktywnych i~warunków ukrytych. Z~tej
perspektywy skuteczność DLP zależy od stanu pozostałych warstw: klasyfikacji
danych, zakresu uprawnień, segmentacji sieci i~procedur obsługi alarmów.
Przypadek Target pokazuje, że zabezpieczenie techniczne może zadziałać
prawidłowo, a~mimo to nie zapobiec szkodzie, gdy nikt nie reaguje na jego
sygnały\footcite{cheng2017enterprise}.

Z~przeprowadzonej analizy wynikają trzy wnioski praktyczne:
\begin{enumerate}
  \item Projektowanie ochrony powinno zaczynać się od identyfikacji
    i~klasyfikacji danych, a~dopiero potem obejmować dobór technologii.
  \item Wdrożenie powinno przebiegać etapami, od obserwacji przez
    dostrajanie reguł do blokowania. Każdy alarm i~każdy wyjątek powinien
    mieć przypisanego właściciela.
  \item Ocena skuteczności powinna obejmować trafność alarmów, czas reakcji
    i~zakres pokrycia chronionych zasobów, a~nie wyłącznie liczbę
    zablokowanych operacji.
\end{enumerate}

Wnioski te mają charakter koncepcyjny. Wymagają weryfikacji w~badaniach
empirycznych prowadzonych w~organizacjach o~różnej wielkości i~strukturze.

Kierunki dalszych badań obejmują przede wszystkim wykrywanie wycieków
w~kanałach szyfrowanych, zastosowanie uczenia głębokiego do wykrywania
zagrożeń wewnętrznych oraz DLP świadczone jako usługa chmurowa z~zachowaniem
poufności analizowanych danych\footcite{cheng2017enterprise}. Potrzebne są
też wspólne zbiory danych testowych, bez których trudno rzetelnie porównywać
metody wykrywania. Osobnym zagadnieniem jest przepływ informacji do narzędzi
generatywnej sztucznej inteligencji, którego skala według danych z~2025~roku
szybko rośnie\footcite{verizon2026dbir}.

% ---------------------------------------------------------------------
%  Bibliografia (alfabetycznie, z pliku bibliografia.bib)
% ---------------------------------------------------------------------
\printbibliography

% ---------------------------------------------------------------------
\section*{Nota biograficzna}

{\footnotesize\setstretch{1}
dr inż. \textbf{Bartosz Lewandowski} - Dyrektor Działu Innowacji i Rozwoju w~Atende S.A. 
oraz adiunkt w~Katedrze Inżynierii Klinicznej Akademii Śląskiej. 
Zajmuje się sztuczną inteligencją, bezpieczeństwem i~suwerennością danych oraz innowacyjnymi 
technologiami medycznymi i~finansowymi. 
Posiada ponad 15-letnie doświadczenie w~projektowaniu i~wdrażaniu
systemów informatycznych oraz zarządzaniu zespołami
badawczo-rozwojowymi. Rozwija rozwiązania wykorzystujące lokalne
modele językowe, architektury RAG i~agentów AI, ze~szczególnym
uwzględnieniem cyberbezpieczeństwa, ochrony informacji
i~suwerenności danych. Prowadzi badania nad zastosowaniami AI w~medycynie, m.in. w~analizie obrazów diagnostycznych oraz predykcji,
a~także nad uwierzytelnianiem bezpieczeństwa żywności metodami uczenia maszynowego. 
Współtworzył rozwiązania z~zakresu bezpiecznego cyfrowego zarządzania dokumentami, 
w~tym system ChainDoc oparty na technologii blockchain, realizujący wymogi trwałego nośnika w~bankowości. 
Zespół ChainDoc otrzymał nagrodę ComputerWorld za najlepsze wdrożenie technologii blockchain 
oraz nagrodę Smart City Forum w~kategorii usług wspólnych. 
Organizator \enquote{Konferencji PI} w Toruniu, prelegent na konferencjach krajowych oraz międzynarodowych.

\par}

\end{document}
